\documentclass[../../../include/open-logic-section]{subfiles}

\begin{document}
\olfileid{sth}{ord-arithmetic}{intro}

\olsection{परिचय}

\olref[ordinals][]{chap} मध्ये आपण क्रमसूचक संख्यांची उपपत्ती विकसित केली. \olref[spine][]{chap} मध्ये आपण पाहिले की पदानुक्रमाचा उर्वरित भाग ज्याभोवती बांधला जातो असा कणा म्हणून क्रमसूचक संख्यांकडे पाहता येते. पण क्रमसूचक संख्यांची तेवढीच भूमिका नाही. क्रमसूचक अंकगणित करण्याचे कामही आहे.

\olref[ordinals][idea]{sec} मध्ये $\omega$, $\omega+1$ आणि $\omega+\omega$ यांचा उल्लेख करताना आपण याची आधीच थोडी कल्पना दिली होती. तेव्हा आपण अनौपचारिकपणे बोललो; आता ते नीट स्पष्ट करण्याची वेळ आली आहे. मात्र या प्रकरणात तत्त्वज्ञान फारसे नाही, हे सांगायला हवे. येथे मुख्यतः तांत्रिक विकास आहे; त्याबरोबर एकच गोष्ट दोन वेगवेगळ्या मार्गांनी करता येते हे दाखवणारे (किंचित) रोचक निरीक्षण आहे.

\end{document}
